\subsection{Multi-Agent Debate Orchestrator (Implementation)}

We implemented the multi-agent debate system in Python using LangGraph (\cite{langchain}), totaling approximately 2{,}900 lines across seven modules. The system comprises six components: typed data models, a configuration layer, role-specific prompts, the LangGraph debate graph, a high-level runner, and a test suite.

\textbf{Data Models.}
We define Pydantic models that mirror the TypeScript types in the baseline single-agent system (\texttt{agents/types.ts}) to ensure cross-language compatibility. The core types are \texttt{Observation} (market state, portfolio state, text context, constraints), \texttt{Action} (orders, justification, confidence, causal claims), and \texttt{AgentTrace} (full auditable record of the debate). Each causal claim carries a \texttt{PearlLevel} tag (L1: association, L2: intervention, L3: counterfactual) for downstream reasoning evaluation.

\textbf{Agent Roles.}
The system supports six specialist roles, each with a dedicated system prompt encoding domain expertise and step-by-step reasoning instructions:
\begin{itemize}
\setlength\itemsep{0pt}
  \item \textbf{Macro Strategist}: Fed policy, inflation, yield curve, global macro.
  \item \textbf{Value Analyst}: Earnings quality, valuation multiples, balance sheet health.
  \item \textbf{Risk Manager}: Volatility regimes, drawdown risk, position sizing, tail risk.
  \item \textbf{Technical Analyst}: Price action, momentum indicators, support/resistance.
  \item \textbf{Sentiment Analyst}: News sentiment, market psychology, narrative shifts.
  \item \textbf{Devil's Advocate}: Stress-tests the emerging consensus and prevents groupthink.
\end{itemize}

Every role prompt includes three anti-failure-mode injections: (1)~a \emph{Causal Claim Format} block requiring at least one L2 or L3 claim with explicit variables, assumptions, and falsifiers; (2)~a \emph{Forced Uncertainty} block requiring disclosure of what evidence would change the agent's mind, the probability the agent is wrong, and potential confounders; and (3)~a \emph{Trap Awareness} block listing common reasoning errors (reverse causation, confounding, selection bias, post-hoc narrative, survivorship bias).

\textbf{Debate Graph.}
The orchestrator is a LangGraph \texttt{StateGraph} with eight node functions connected by directed edges:

\begin{center}
\texttt{START} $\to$ \texttt{news\_digest} / \texttt{data\_analysis} $\to$ \texttt{build\_context} $\to$ \texttt{propose} $\to$ \texttt{critique} $\rightleftharpoons$ \texttt{revise} $\to$ \texttt{judge} $\to$ \texttt{build\_trace} $\to$ \texttt{END}
\end{center}

The \texttt{news\_digest} and \texttt{data\_analysis} pipeline nodes are optional and run in parallel when enabled; they preprocess raw text and numerical data into structured intelligence briefs before the debate begins. In the \texttt{propose} phase, each agent independently generates a trading proposal with orders, causal claims, and falsifiers. In the \texttt{critique} phase, each agent critiques every other agent's proposal, identifies the weakest assumption, and performs self-critique. In the \texttt{revise} phase, agents update their proposals in response to the critiques they received, optionally upgrading L1 claims to L2/L3. A conditional edge on the \texttt{revise} node loops back to \texttt{critique} for additional rounds (configurable via \texttt{max\_rounds}), or proceeds to the \texttt{judge} node, which synthesizes all revised proposals into a single final decision with an audited memo. The judge explicitly preserves the strongest objection from the debate, even if overruled, to maintain disagreement visibility. The \texttt{build\_trace} node packages the full debate history into a JSON trace for post-hoc evaluation.

\textbf{Experimental Knobs.}
All experimental parameters are centralized in a \texttt{DebateConfig} dataclass to support systematic ablation. The key parameters are:
\begin{itemize}
\setlength\itemsep{0pt}
  \item \textbf{Agent roster}: any subset of the six roles (default: macro, value, risk, technical).
  \item \textbf{Max rounds}: number of critique$\to$revise cycles (default: 1).
  \item \textbf{Agreeableness}: a continuous knob from 0.0 (maximally confrontational) to 1.0 (sycophantic), injected into critique and revision prompts. At low values, agents are instructed to challenge every assumption; at high values, they seek consensus. The default is 0.3 (skeptical).
  \item \textbf{Adversarial mode}: when enabled, a Devil's Advocate agent is automatically added to the roster.
  \item \textbf{Pipeline toggles}: independently enable or disable news digest and data analysis preprocessing.
  \item \textbf{LLM settings}: model name (default \texttt{gpt-4o-mini}) and temperature (default 0.3).
  \item \textbf{Mock mode}: deterministic synthetic responses for testing without API calls.
\end{itemize}

\textbf{LLM Integration.}
In live mode, each node calls OpenAI's API via \texttt{langchain-openai} with retry logic (up to 3 attempts with exponential backoff on transient errors). In mock mode, deterministic generator functions produce role-biased synthetic proposals, critiques, and revisions, enabling the full test suite to run without API access.
